\documentclass[11pt]{article}
\usepackage[margin=1in]{geometry}
\usepackage{amsmath,amssymb,amsthm,mathtools}
\usepackage{booktabs}
\usepackage{verbatim}
\usepackage{multicol}
\usepackage[colorlinks=true,linkcolor=blue,citecolor=blue,urlcolor=blue]{hyperref}
\hypersetup{pdftitle={167 XORs for Rank-47 4x4 Matrix Multiplication over F2},pdfauthor={Greg Sidebottom}}
\newtheorem{theorem}{Theorem}[section]
\newtheorem{lemma}[theorem]{Lemma}
\newtheorem{proposition}[theorem]{Proposition}
\theoremstyle{remark}
\newtheorem{remark}[theorem]{Remark}
\newcommand{\F}{\mathbb{F}}
\newcommand{\Z}{\mathbb{Z}}
\newcommand{\XOR}{\operatorname{X}}
\title{167 XORs for Rank-47 $4\times4$ Matrix Multiplication\\
in Characteristic Two}
\author{Greg Sidebottom}
\date{Draft --- October 17, 2026}

\begin{document}
\maketitle

\begin{abstract}
Rank-47 bilinear schemes multiply $4\times4$ matrices with $47$
multiplications over any ring of characteristic two; at least three
inequivalent ones are publicly available.  Applied recursively, a
scheme whose three linear maps use $q$ additions (XORs) costs
$(1+q/31)\,n^{\log_4 47}+o(n^{\log_4 47})$ operations, so $q$ fixes
the leading constant.  The best count known for a rank-47 scheme was
$170$ (Dumas, Pernet, Sedoglavic, Stapleton and Tichavsk\'y, 2026).  We
give a rank-47 scheme over $\F_2$ whose maps cost
$44+49+74=167$ XORs, leading constant $6+12/31$, and $168$ for the
published scheme itself.  The new scheme lies in the symmetry orbit of
Kauers and Moosbauer's rank-47 scheme.  It was found by simulated
annealing over the orbit, scored by a Boyar--Peralta heuristic with
exact distances, followed by an exact search for the minimum number of
non-target XORs on each side; for this scheme all three counts are
proved minimal.  Each side depends on only two of the three basis
changes and is unchanged by permutation matrices, so a class's whole
orbit reduces to three $840\times840$ cost tables and a min-plus scan
over $840^3$ triples.  The scan shows that the heuristic minimum is
exactly $168$ in all three known classes; exact refinement of the
near-minimum triples reaches $167$ only in the Kauers--Moosbauer class.
All programs are checked
by an independent verifier and by PLinOpt's \texttt{SLPchecker}.  We
also summarize the current rank and addition records for $4\times4$
matrix multiplication with and without change of basis.
\end{abstract}

\section{Introduction}

Strassen's algorithm applied to $2\times2$ blocks of $2\times2$
matrices multiplies $4\times4$ matrices with $49$ multiplications over
any ring~\cite{Strassen69}.  AlphaTensor lowered this to $47$ in
characteristic two~\cite{Fawzi22}; Kauers and Moosbauer found an
inequivalent rank-47 scheme with flip graphs~\cite{KM23}; and a third,
attributed in~\cite{DPSST26} to AlphaEvolve~\cite{Novikov25}, is the
one whose programs PLinOpt distributes~\cite{PLinOpt}.  Outside
characteristic two the record is $48$, over the complex numbers by
AlphaEvolve and over any ring containing $1/2$ by Dumas, Pernet and
Sedoglavic~\cite{DPS25}.

Used recursively, a scheme's multiplication count fixes the exponent
and its additions fix the constant.  For an
$\langle n_0\times n_0\times n_0:r\rangle$ scheme whose linear maps use
$q$ additions and scalings in total,
\[
F(n) = r\,F(n/n_0) + q\,(n/n_0)^2
\quad\Longrightarrow\quad
F(n) = \Bigl(1+\frac{q}{r-n_0^2}\Bigr)n^{\log_{n_0} r} + o\bigl(n^{\log_{n_0} r}\bigr),
\]
so at $r=47$, $n_0=4$ the constant is $1+q/31$.  Over $\F_2$ every
addition is an XOR and no scalings occur.  Recent work drove the
rank-48 count from $315$ operations~\cite{DPS25} down to
$204$~\cite{DPSST26}; for rank $47$, the same paper reports the best
known program at $170$ XORs, constant $6+15/31$.

\paragraph{Contributions.}
\begin{itemize}
\item A rank-47 scheme over $\F_2$ with $167=44+49+74$ XORs (constant
  $6+12/31\approx6.387$), given in full in Appendix~A, and $168$ for
  the published scheme (Section~\ref{sec:results}).  For the new
  scheme each side's count is proved minimal.
\item An exact method for the minimum XOR count of a side with few
  non-target signals (Section~\ref{sec:exact}): the problem reduces to
  choosing those signals, and iterative deepening over their number is
  exact and fast up to five of them.
\item An exhaustive scan of each scheme's symmetry orbit, the $4\times4$
  form of the orbit scan used for $3\times3$ in~\cite{Sidebottom26a},
  showing that all three known classes bottom out at $168$ under
  heuristic scoring (Section~\ref{sec:results}).
\item A summary of the current $4\times4$ records
  (Section~\ref{sec:records}).
\end{itemize}

\section{Records for $4\times4$ matrix multiplication}\label{sec:records}

Table~\ref{tab:rank} lists the best known ranks; Tables~\ref{tab:flat}
and~\ref{tab:alt} the best known addition counts, without and with a
change of basis.  The leading constants compare only within one rank,
since each rank recurses with its own exponent.

\begin{table}[ht]
\centering
\caption{Best known rank of $\langle4,4,4\rangle$.}\label{tab:rank}
\begin{tabular}{llp{7.6cm}}
\toprule
setting & rank & source \\
\midrule
characteristic $2$ (e.g.\ $\F_2$) & $47$ & AlphaTensor~\cite{Fawzi22}; inequivalent
  schemes by Kauers--Moosbauer~\cite{KM23} and AlphaEvolve~\cite{Novikov25,DPSST26} \\
rings containing $1/2$ & $48$ & AlphaEvolve (over $\mathbb{C}$)~\cite{Novikov25};
  rational coefficients~\cite{DPS25} \\
all rings & $49$ & Strassen applied twice~\cite{Strassen69} \\
\bottomrule
\end{tabular}
\end{table}

\begin{table}[ht]
\centering
\caption{Best known additions (with scalings) without change of basis;
constant of $n^{\log_4 r}$.}\label{tab:flat}
\small
\begin{tabular}{lllll}
\toprule
rank (valid over) & exponent & ops & constant & source \\
\midrule
$47$ (char.\ $2$) & $\log_4 47\approx2.777$ & $170$ & $6+15/31$ & \cite{PLinOpt}, cited in~\cite{DPSST26} \\
                  &                       & $\mathbf{167}$ & $\mathbf{6+12/31}$ & this paper \\
$48$ (rings with $1/2$) & $\log_4 48\approx2.792$ & $315$ & $10+27/32$ & \cite{DPS25} \\
                  &                       & $204$ & $7+3/8$ & \cite{DPSST26}\,$^{a}$ \\
$49$ (all rings)  & $\log_2 7\approx2.807$ & $165$ & $6$ & \cite{Winograd71}, applied twice \\
                  &                       & $156$ & $5+8/11$ & \cite{PLinOpt}, cited in~\cite{DPSST26} \\
\bottomrule
\end{tabular}

\smallskip
\footnotesize $^{a}$ The PLinOpt repository also holds a $202$-operation
variant (data file \texttt{4x4x4\_48\_202}, contributed by R.~Patel,
October~2026).
\end{table}

\begin{table}[ht]
\centering
\caption{Best known constants with a change of basis (conversion
costs are lower order).  $d$ is the encoding dimension; $d=16$ is a
square basis change.}\label{tab:alt}
\small
\begin{tabular}{lllll}
\toprule
rank & basis change & recursive ops & constant & source \\
\midrule
$49$ & square ($2\times2$, twice) & $12$ per $2\times2$ step & $5$ & \cite{KS17} \\
$48$ & square ($d=16$) & $173$ & $6+13/32$ & \cite{DPSST26} \\
$48$ & rectangular ($d=24$) & $122$ & $6+1/12$ & \cite{DPSST26} \\
$48$ & rect.\ ($d=24$: $A,B$; $32$: $C$) & $32+32+32$ & $5+2/3$ & \cite{DPSST26} \\
$47$ & --- & --- & --- & none published \\
\bottomrule
\end{tabular}
\end{table}

Two remarks.  First, the rectangular encodings of Beniamini and
Schwartz~\cite{BS19,BCHKS20} are a different cost model from
Karstadt and Schwartz's square basis change~\cite{KS17}: with them a
$3\times3$ rank-23 base already reaches constant $2$~\cite{BS19}, the
constant of the classical algorithm.  Second, the $3\times3$ analogue
of this paper's question has seen a run of improvements:
$60$~\cite{Stapleton25}, $59$~\cite{MSWS26}, $58$~\cite{Perminov25a},
$56$~\cite{Sun26} and
$55$ additions~\cite{Sidebottom26a,KI26} without change of basis, and
$52$~\cite{Nielsen26} then $51$ additions per step~\cite{SP26,Sidebottom26b}
with a square basis change.  M{\aa}rtensson and Stankovski
Wagner~\cite{MW24,MW26} and Perminov~\cite{Perminov25b} study
addition reduction for many formats; Dumas, Pernet and
Sedoglavic~\cite{DPS25,DPS26a,DPSST26} and the PLinOpt
tools~\cite{PLinOpt} are the reference for $4\times4$.

\section{Model, lower bounds and symmetry}\label{sec:model}

A rank-$r$ scheme for $\langle4,4,4\rangle$ over $\F_2$ is a triple of
matrices $L\in\F_2^{r\times16}$, $R\in\F_2^{r\times16}$,
$P\in\F_2^{16\times r}$ with
\[
\operatorname{vec}(AB) = P\bigl((L\operatorname{vec}A)\odot(R\operatorname{vec}B)\bigr)
\]
for all $A,B$, where $\operatorname{vec}$ lists entries row by row and
$\odot$ is the entrywise product.  Equivalently, product $k$ is
$m_k=\langle U_k,A\rangle\langle V_k,B\rangle$ for $4\times4$ matrices
$U_k,V_k$ (rows $k$ of $L$ and $R$) and $AB=\sum_k m_kW_k$ ($W_k$ is
column $k$ of $P$).  Since the identity holds over the field $\F_2$, it
holds for entries in any ring of characteristic two, commutative or not,
so the scheme recurses on blocks.

\paragraph{Cost.}  For a matrix $M$ with rows over the inputs, let
$\XOR(M)$ be the fewest XOR gates of a straight-line program computing
all rows of $M$ from the inputs (copies are free).  The cost of a scheme
is $\XOR(L)+\XOR(R)+\XOR(P)$.  The output side is handled through its
transpose: by the transposition principle~\cite{BCS97} a linear program
for a matrix without zero rows or columns can be reversed into one for
its transpose, and $\XOR(P)=\XOR(P^{\mathsf T})+(r-16)$.  All three
sides are therefore instances of one problem: compute a set of $r$
vectors of $\F_2^{16}$ from the $16$ unit vectors.  Over $\F_2$ the
shortest such program is NP-hard to find in general~\cite{BMP08}.

\paragraph{Floor.}  Every target that is not a unit vector needs a gate
of its own, so $\XOR(M)\ge|T(M)|$, where $T(M)$ is the set of distinct
non-unit rows.  For a scheme this gives the floor
$|T(L)|+|T(R)|+|T(P^{\mathsf T})|+31$.

\paragraph{Symmetry.}  If $(U_k,V_k,W_k)_k$ is a scheme, so is
$(X^{\mathsf T}U_kY^{-\mathsf T},\;Y^{\mathsf T}V_kZ^{-\mathsf T},\;
X^{-1}W_kZ)_k$ for invertible $X,Y,Z$ (substitute $A\mapsto XAY^{-1}$,
$B\mapsto YBZ^{-1}$), and so are the images under the cyclic symmetry
$(U,V,W)\mapsto(V,W^{\mathsf T},U^{\mathsf T})$ and the transposition
$(U,V,W)\mapsto(V^{\mathsf T},U^{\mathsf T},W^{\mathsf T})$ of the
matrix multiplication tensor~\cite{deGroote78}.  Over $\F_2$ the orbit
of a scheme has up to $6\cdot20160^3\approx4.9\cdot10^{13}$ points, all
with the same multiplications but different linear maps, and hence
different XOR counts.  The multiset of rank triples
$(\operatorname{rank}U_k,\operatorname{rank}V_k,\operatorname{rank}W_k)$,
taken up to a common permutation of the three coordinates, is constant
on an orbit; it separates the three published rank-47 schemes, which
are therefore pairwise inequivalent.  None of the three has two
products sharing a factor, so each is an isolated vertex of the flip
graph~\cite{KM23} and flip walks cannot leave it.

\section{Method}\label{sec:method}

\subsection{Boyar--Peralta with exact distances}

Boyar and Peralta's heuristic~\cite{BP10,BMP13} keeps a base $S$ of
computed signals, initially the unit vectors, and the distance
$\delta(t)$ of each target: one less than the fewest elements of $S$
summing to $t$.  Each step adds a target at distance one if there is
one, and otherwise the sum of two base elements that minimizes the
total distance, ties broken by the larger sum of squared distances.
With $16$ inputs the distances are exact and cheap: a table of the
$2^{16}$ values $k(v)$ (fewest base elements summing to $v$) starts at
the Hamming weights and is updated in one pass per new signal $s$ by
$k(v)\leftarrow\min(k(v),\,1+k(v+s))$.  We add random tie breaking and
optional slack, in the spirit of Tan and Peyrin~\cite{TP20}, run many
restarts, delete unused gates, and reverse the program for $P^{\mathsf T}$
into one for $P$.  A pilot variant (complete greedily after each of the
best few candidates and commit to the best completion) never improved
on the plain heuristic in our instances.

\subsection{Searching the symmetry orbit}

A point of the orbit is a slot symmetry together with $(X,Y,Z)\in
GL_4(\F_2)^3$; $X$ affects the $L$ and $P$ sides, $Y$ the $L$ and $R$
sides, $Z$ the $R$ and $P$ sides.  We tried coordinate descent (each
step scans all $20160$ values of one coordinate, shortlists the best by
floor and weight, and scores the shortlist with the heuristic) and
simulated annealing whose moves multiply one coordinate by an
elementary matrix $I+E_{ij}$ on either side and rescore the two
affected sides.  Annealing was clearly better: on the AlphaTensor class
two chains of $3000$ moves reached $168$, where $240$ coordinate-descent
restarts reached $169$ once.  Every orbit point that is written out is
checked against all $256$ basis products.

\paragraph{Exhaustive scan.}  The orbit can also be covered completely.
Multiplying $X$ on the right by a permutation matrix $Q$ permutes the
rows of every $L$ target ($X^{\mathsf T}U_kY^{-\mathsf T}\mapsto
Q^{\mathsf T}X^{\mathsf T}U_kY^{-\mathsf T}$) and of every $P^{\mathsf T}$
target, a relabelling of the inputs that leaves both XOR counts
unchanged; the same holds for $Y$ and $Z$.  So $X,Y,Z$ range over the
$20160/24=840$ right cosets of the permutation matrices, the three side
costs fill three $840\times840$ tables (about two million heuristic
evaluations), and a min-plus pass over all $840^3\approx5.9\cdot10^8$
triples finds the orbit minimum and every near-minimal triple.  The slot
symmetries add nothing: they permute and transpose the three factor
families, and transposing the targets is again a relabelling, so their
costs are those of points in the untransformed orbit.  This is the
$4\times4$ form of the $168^3$ scan used for $3\times3$
in~\cite{Sidebottom26a}; one class takes about $25$ minutes on six
cores.

\subsection{Exact minimization with few helper signals}\label{sec:exact}

Call the gates of a program that compute targets \emph{target gates}
and the others \emph{helpers}.  A program for $T$ has $|T|+h$ gates
when it uses $h$ helpers.  The following lemma removes all choices
except the helpers.  For $S\subseteq\F_2^{16}$ let $\operatorname{cl}(S)$
be the least superset of $S$ that contains every target equal to the
sum of two of its elements.

\begin{lemma}\label{lem:closure}
A program computing $T$ with exactly $h$ helpers exists if and only if
there are $x_1,\dots,x_h$ such that, with
$S_0=\operatorname{cl}(\{e_1,\dots,e_{16}\})$ and
$S_i=\operatorname{cl}(S_{i-1}\cup\{x_i\})$, each $x_i$ is a sum of two
elements of $S_{i-1}$ and $T\subseteq S_h$.
\end{lemma}

\begin{proof}
Given such $x_i$, build the program by adding targets as soon as they
are sums of two available signals and the $x_i$ in order: every gate is
a target or one of the $h$ helpers.  Conversely, take a program with $h$
helpers, listed in program order as $x_1,\dots,x_h$.  By induction on
the position in the program, every signal before $x_i$ lies in
$S_{i-1}$: a target gate is the sum of two earlier signals, which lie
in the current closure, so the closure contains it; hence each $x_i$ is
a sum of two elements of $S_{i-1}$, and finally $T\subseteq S_h$.
\end{proof}

Two observations prune the search over helper sequences without losing
solutions.  The last helper must make some missing target a sum of two
available signals, so $x_h\in(T\setminus S_{h-1})+S_{h-1}$.  And if
$x_{i+1}$ is already a sum of two elements of $S_{i-1}$, the two
helpers can be exchanged without changing $S_{i+1}$; bubble-sorting any
solution therefore yields one in which such adjacent pairs are
increasing (as integers), and only those are explored.  Iterative
deepening over $h$ returns the exact minimum.  Candidates are tried in
order of how many targets they make immediately available.  The
implementation updates the closure incrementally (adding one signal can
only create targets that use it as an operand), enumerates each helper
set once in its smallest-available-first order, and undoes changes in
place.  With $16$ inputs and $37$--$45$ targets, $h\le4$ is decided in
seconds and $h=5$ in about twenty seconds on twelve cores; each further
level costs roughly a hundred times more.

\subsection{Pipeline}

The scan lists every coset triple within $2$ of its class's minimum;
each is then refined side by side, searching exactly for a program one
helper shorter than the heuristic's, and lower while that succeeds.
Side results are cached by target set, since neighbouring points share
two of their three sides.  Sides that need more helpers than the search
can reach in its time limit stay open, with the lower bound proved so
far.

\section{Results}\label{sec:results}

\paragraph{The published scheme.}  PLinOpt's programs for its rank-47
scheme cost $44+44+82=170$.  The heuristic alone gives $44+44+80=168$
for the same scheme: $P^{\mathsf T}$ needs $49$ XORs, not $51$.  Here the
floor is low ($37+37+34+31=139$) but each side carries many helpers.
The exact search rules out five or fewer helpers on $L$, so
$\XOR(L)\ge43$ against the heuristic's $44$; deeper levels were out of
reach.

\paragraph{The orbit plateau.}  Table~\ref{tab:classes} summarizes the
exhaustive scans.  In every class the heuristic orbit minimum is exactly
$168$, reached by very few coset triples (each coset triple stands for
$24^3$ orbit points with identical costs).  Annealing agrees: $24$ of
$24$ chains end at $168$ for the AlphaTensor and published classes, $23$
of $24$ for the Kauers--Moosbauer class.  Exact refinement of all
triples within $2$ of the minimum lowers the total only in the
Kauers--Moosbauer class: two triples, one scored $168$ and one $169$,
refine to $167$ with sides $43+49+44$ (the last counted for
$P^{\mathsf T}$), the multiset of the appendix scheme's sides, and the
other $168$-triple stays at $168$ with one side open.  At an
AlphaTensor-class point with sides $44+48+45$ the exact search proves
all three sides optimal, so that scheme costs exactly $168$.  For the
triples that could still reach $166$, the exact search settled every
open side needing at most six helpers below the heuristic: of these
$144$ sides it confirmed the heuristic count on $137$, improved it by
exactly one on $4$ (all in the Kauers--Moosbauer class) and left $3$
unresolved within a one-hour limit.  The heuristic is rarely wrong here,
and then by one.  No triple of any class drops below $167$.  In the AlphaTensor class no
near-minimal triple can reach $166$ at all; the triples still open
(last column of Table~\ref{tab:classes}) each have a side that needs
seven or more helpers, beyond the exact search.

\begin{table}[ht]
\centering
\caption{Exhaustive orbit scans (heuristic side costs) and exact
refinement of the triples within $2$ of the minimum.  \emph{Open}: triples
whose proved lower bound still allows $166$ or less; each has a side
that needs at least seven helpers below the heuristic.}\label{tab:classes}
\small
\begin{tabular}{lcccccc}
\toprule
class & as given & orbit min & triples at min & refined & best exact & open \\
\midrule
published~\cite{PLinOpt} (AlphaEvolve) & $170$ & $168$ & $2$ & $216$ & $168$ & $78$ \\
AlphaTensor~\cite{Fawzi22} & $179$ & $168$ & $13$ & $195$ & $168$ & $0$ \\
Kauers--Moosbauer~\cite{KM23} & $186$ & $168$ & $2$ & $100$ & $\mathbf{167}$ & $7$ \\
\bottomrule
\end{tabular}
\end{table}

\paragraph{The 167 scheme.}  Exact refinement lowers the most expensive
side of many Kauers--Moosbauer-class points by one: of the first $537$
points refined, $211$ reach $167$, $234$ stay at $168$ and $92$ at $169$;
none reaches $166$.  The scheme of Appendix~A is one of them.  It is the
image of Kauers and Moosbauer's scheme under the slot symmetry
$(U,V,W)\mapsto(W,V^{\mathsf T},U)$ followed by the substitution with
$(X,Y,Z)$ given by \texttt{54a8}, \texttt{59f8}, \texttt{a512} in
Appendix~A's encoding (recomputing this image reproduces the appendix
exactly).  Its sides have $41$, $45$ and
$41$ distinct non-unit targets (floor $158$), and the exact search
gives
\[
\XOR(L)=44\ (3\text{ helpers}),\qquad
\XOR(R)=49\ (4\text{ helpers}),\qquad
\XOR(P^{\mathsf T})=43\ (2\text{ helpers}),
\]
each with the next smaller helper count refuted, so
$\XOR(P)=43+31=74$ and $167$ is the exact minimum for this scheme
without change of basis.  The heuristic finds $50$ for $R$; the exact
search finds the four-helper program.

\paragraph{Leading constant.}  The scheme gives
$F(n)=\tfrac{198}{31}\,n^{\log_4 47}+o(n^{\log_4 47})$ ring
operations (additions and multiplications) over rings of characteristic
two, against $\tfrac{201}{31}$ before.

\section{Verification and artifacts}

The scheme and programs are in the repository
\url{https://github.com/gsidebottom/logic}, directory
\texttt{matmul/x47/} (\texttt{best/x167/}: \texttt{x167\_\{L,R,P\}.sms}
and \texttt{.slp} in PLinOpt formats).  Two independent checks pass:
\begin{itemize}
\item \texttt{x47check.py} (written for this purpose, independent of
  the search code) evaluates each program over $\F_2$, compares every
  output with the corresponding matrix row, counts XORs as the number of
  terms minus one per line, and checks the scheme on all $256$ pairs of
  basis matrices, which by bilinearity proves
  $P((L\operatorname{vec}A)\odot(R\operatorname{vec}B))=\operatorname{vec}(AB)$;
  it reports \texttt{OK rank 47: L 44 + R 49 + P 74 = 167 XORs
  (Brent 256/256 basis pairs)}.
\item PLinOpt's \texttt{SLPchecker -q 2 -M x167\_S.sms < x167\_S.slp}
  reports \texttt{<44, 0>}, \texttt{<49, 0>} and \texttt{<74, 0>}
  for $S=L,R,P$.
\end{itemize}
The search tool (\texttt{src/bin/xor47.rs}: \texttt{bp}, \texttt{scheme},
\texttt{orbit}, \texttt{scan}, \texttt{exact}, \texttt{refine},
\texttt{settle}) regenerates every number in
Section~\ref{sec:results}; the symmetric SAT instances come from
\texttt{matmul/x47/s3cnf.py}.

\section{Outlook}

The three known classes all bottom out at $168$ under the heuristic,
and the best refined point is $167$, so further progress most likely
needs new rank-47 classes.  All known classes are flip-isolated, which
rules out plain flip walks.  SAT-based search found thousands of
$3\times3$ schemes~\cite{HKS19}; we tried SAT on schemes invariant under
the slot symmetries.  Invariance under the cyclic symmetry alone divides
the unknowns by three, but no solver found a rank-49 cyclic scheme
(which exists: the square of a cyclic Strassen scheme) in two hours.
Adding the transposition symmetry ($S_3$-invariance, products in orbits
of size $6$, $3$, $2$, $1$) brings the rank-49 instance to $410$ unknowns,
and CaDiCaL solves it in six minutes; at rank $47$, eight hours over the
$144$ orbit shapes found no scheme, with $53$ shapes refuted by the
solver (unverified) and $91$ open.  Other directions are an exact search
reaching six or more helpers (it would settle, for example, whether the
published scheme's $L$ side admits $43$), the first alternative-basis
constants for rank $47$, and rectangular encodings in the sense
of~\cite{BS19}.

\section*{Acknowledgments}

This work was carried out in an extended interactive collaboration
with Claude (Anthropic), which implemented the search and verification
tools and much of this text under the author's direction and review.
Every claim is checked mechanically by the tools above; the results
rest on those checks rather than on trust in either the author or the
tools.

\begin{thebibliography}{99}
\bibitem{BCHKS20} G.~Beniamini, N.~Cheng, O.~Holtz, E.~Karstadt,
O.~Schwartz, \emph{Sparsifying the operators of fast matrix
multiplication algorithms}, arXiv:2008.03759 (2020).
\bibitem{BS19} G.~Beniamini, O.~Schwartz, \emph{Fast matrix
multiplication via sparse decomposition}, Proc.\ SPAA 2019, 11--22.
\bibitem{BMP08} J.~Boyar, P.~Matthews, R.~Peralta, \emph{On the
shortest linear straight-line program for computing linear forms},
Proc.\ MFCS 2008, LNCS 5162, 168--179.
\bibitem{BMP13} J.~Boyar, P.~Matthews, R.~Peralta, \emph{Logic
minimization techniques with applications to cryptology}, J.\
Cryptology 26 (2013) 280--312.
\bibitem{BP10} J.~Boyar, R.~Peralta, \emph{A new combinational logic
minimization technique with applications to cryptology}, Proc.\ SEA
2010, LNCS 6049, 178--189.
\bibitem{BCS97} P.~B\"urgisser, M.~Clausen, M.~A.~Shokrollahi,
\emph{Algebraic Complexity Theory}, Springer, 1997.
\bibitem{deGroote78} H.~F.~de~Groote, \emph{On varieties of optimal
algorithms for the computation of bilinear mappings I. The isotropy
group of a bilinear mapping}, Theoret.\ Comput.\ Sci.\ 7 (1978) 1--24.
\bibitem{DPS25} J.-G.~Dumas, C.~Pernet, A.~Sedoglavic, \emph{A
non-commutative algorithm for multiplying $4\times4$ matrices using
$48$ non-complex multiplications}, arXiv:2506.13242 (2025).
\bibitem{DPS26a} J.-G.~Dumas, C.~Pernet, A.~Sedoglavic, \emph{A more
accurate rational non-commutative algorithm for multiplying $4\times4$
matrices using $48$ multiplications}, arXiv:2603.18699 (2026).
\bibitem{DPSST26} J.-G.~Dumas, C.~Pernet, A.~Sedoglavic, J.~Stapleton,
P.~Tichavsk\'y, \emph{Fast matrix multiplication via recursive
$\langle4\times4\times4:48\rangle$ algorithms into practice},
arXiv:2609.12027 (2026).
\bibitem{PLinOpt} J.-G.~Dumas, B.~Grenet, C.~Pernet, A.~Sedoglavic,
\emph{PLinOpt, a collection of C++ routines handling linear \&
bilinear programs}, \url{https://github.com/jgdumas/plinopt} (data
file \texttt{4x4x4\_47\_mod\_2}, September 2026).
\bibitem{Fawzi22} A.~Fawzi et al., \emph{Discovering faster matrix
multiplication algorithms with reinforcement learning}, Nature 610
(2022) 47--53.
\bibitem{HKS19} M.~J.~H.~Heule, M.~Kauers, M.~Seidl, \emph{Local search
for fast matrix multiplication}, Proc.\ SAT 2019, LNCS 11628, 155--163.
\bibitem{KS17} E.~Karstadt, O.~Schwartz, \emph{Matrix multiplication, a
little faster}, Proc.\ SPAA 2017, 101--110; J.\ ACM 67 (2020).
\bibitem{KI26} S.~Karunaratne, A.~Idamekorala, \emph{55 additions
suffice for $3\times3$ matrix multiplication at rank 23},
arXiv:2607.28676 (2026).
\bibitem{KM23} M.~Kauers, J.~Moosbauer, \emph{Flip graphs for matrix
multiplication}, Proc.\ ISSAC 2023; arXiv:2212.01175.
\bibitem{MSWS26} E.~M{\aa}rtensson, P.~Stankovski Wagner, J.~Stapleton,
\emph{A rank 23 algorithm for multiplying $3\times3$ matrices with an
arithmetic complexity of 59}, arXiv:2601.05272 (2026).
\bibitem{MW24} E.~M{\aa}rtensson, P.~Stankovski Wagner, \emph{The number
of the beast: reducing additions in fast matrix multiplication
algorithms for dimensions up to 666}, IACR ePrint 2024/2063.
\bibitem{MW26} E.~M{\aa}rtensson, P.~Stankovski Wagner, \emph{On why and
how to minimize the arithmetic complexity of fast matrix multiplication
algorithms}, IACR ePrint 2026/849.
\bibitem{Nielsen26} H.~Nielsen, \emph{A 52-addition, rank-23 scheme for
exact $3\times3$ matrix multiplication}, arXiv:2609.06588 (2026).
\bibitem{Novikov25} A.~Novikov et al., \emph{AlphaEvolve: a coding agent
for scientific and algorithmic discovery}, arXiv:2506.13131 (2025).
\bibitem{Perminov25a} A.~I.~Perminov, \emph{A 58-addition, rank-23
scheme for general $3\times3$ matrix multiplication}, arXiv:2512.21980
(2025).
\bibitem{Perminov25b} A.~I.~Perminov, \emph{Parallel heuristic
exploration for additive complexity reduction in fast matrix
multiplication}, arXiv:2512.13365 (2025).
\bibitem{Sidebottom26a} G.~Sidebottom, \emph{A 55-addition rank-23
scheme for $3\times3$ matrix multiplication via exact two-sided
minimization}, Zenodo,
\href{https://doi.org/10.5281/zenodo.21240904}{doi:10.5281/zenodo.21240904}
(July 2026).
\bibitem{Sidebottom26b} G.~Sidebottom, \emph{51 additions per recursive
step for rank-23 $3\times3$ matrix multiplication, machine-checked over
arbitrary rings}, Zenodo,
\href{https://doi.org/10.5281/zenodo.23305910}{doi:10.5281/zenodo.23305910}
(October 2026).
\bibitem{SP26} J.~Stapleton, A.~Perminov, \emph{A 51-addition
alternative-basis kernel for rank-23 $3\times3$ matrix
multiplication}, arXiv:2609.32388 (2026).
\bibitem{Stapleton25} J.~Stapleton, \emph{A 60-addition, rank-23 scheme
for exact $3\times3$ matrix multiplication}, arXiv:2508.03857 (2025).
\bibitem{Strassen69} V.~Strassen, \emph{Gaussian elimination is not
optimal}, Numer.\ Math.\ 13 (1969) 354--356.
\bibitem{Sun26} Y.~Sun, \emph{An exact 56-addition, rank-23 scheme for
general $3\times3$ matrix multiplication}, arXiv:2604.27645 (2026).
\bibitem{TP20} Q.~Q.~Tan, T.~Peyrin, \emph{Improved heuristics for short
linear programs}, IACR Trans.\ Cryptogr.\ Hardw.\ Embed.\ Syst.\
2020(1) 203--230.
\bibitem{Winograd71} S.~Winograd, \emph{On multiplication of $2\times2$
matrices}, Linear Algebra Appl.\ 4 (1971) 381--388.
\end{thebibliography}

\appendix
\section{The 167-XOR scheme}\label{app:scheme}
% generated by matmul/x47/mk_paper_appendix.py -- do not edit
\subsection*{A.1\quad The scheme}
Product $k$ is $m_k=\langle U_k,A\rangle\,\langle V_k,B\rangle$ and $C=\sum_k m_k W_k$ over $\F_2$.  Each $4\times4$ coefficient matrix is written as four hexadecimal digits, digit $i$ ($i=0,\dots,3$, left to right) encoding row $i$ with bit value $2^j$ for entry $(i,j)$; e.g.\ \texttt{1000} is $E_{00}$ and \texttt{8421} the identity.
\begin{center}\footnotesize
\begin{verbatim}
 k  U    V    W        k  U    V    W        k  U    V    W
 1  f07e 0366 2000    17  0808 1001 7e0b    33  0081 3067 7e7e
 2  d05e 0360 3e1e    18  0e7e b06d b2b0    34  0809 1067 090b
 3  0009 0067 0909    19  1010 5005 7eb0    35  585c 0144 7000
 4  7e7e 00dd 8100    20  0909 10cd 0900    36  0080 3063 7e5e
 5  0d7c da6d 7a78    21  0001 7067 7e7c    37  0d0d d0cd 7e78
 6  0002 0760 0404    22  0004 0010 7e0f    38  0d5c 9a6d 7070
 7  7a7e 01dd 0100    23  9090 0055 9000    39  007e 006d 1818
 8  0e0e d00d 3ab0    24  0d7e 0a6d 0808    40  000e 0060 3a1a
 9  787c 0100 7a0b    25  707e 0066 2b0b    41  1e5e d80d 3eb0
10  0070 d06d 3a38    26  0d4d 90cd 7e70    42  0090 306d 9090
11  787e 0166 0b0b    27  1c5c 984d 0e00    43  1090 3065 90b0
12  0f7e ba6d 0200    28  1e7e b86d b0b0    44  d0d4 0054 7e0e
13  1818 1045 7e00    29  3e7e 08dd 8000    45  0202 dc0d 0400
14  0040 d08d 7ef0    30  2020 0500 4000    46  d05c 0364 0e0e
15  d090 005d 9010    31  0d09 00cd 0908    47  d07e 036d 1010
16  d0d0 0050 7e1e    32  0020 de6d 4040
\end{verbatim}
\end{center}
\subsection*{A.2\quad The programs}
PLinOpt syntax.  In $L$ and $R$ the inputs \texttt{i}$4i{+}j$ are $a_{ij}$ (resp.\ $b_{ij}$) and output \texttt{o}$k$ is the operand of product $k{+}1$; in $P$ the inputs \texttt{i}$k$ are the products $m_{k+1}$ and output \texttt{o}$4i{+}j$ is $c_{ij}$.  A line with $t$ terms costs $t-1$ XORs; copies (one term) are free.
\begin{multicols}{3}\footnotesize
\noindent $L$: 44 XORs
\begin{verbatim}
o2:=i12+i15;
o33:=i7+o2;
o18:=i0+i8;
o41:=i8+i11;
o42:=i11+o18;
o22:=i3+o42;
o29:=i1+i9;
o16:=i7+i15;
o14:=i2+o22;
o15:=i10+o14;
o19:=i4+o33;
o44:=i5+i13;
o43:=i14+o15;
o12:=o18+o16;
o32:=i11+i12;
o30:=i6+o19;
o36:=i14+o30;
o25:=i10+o36;
t18:=o22+o43;
o34:=o12+t18;
o8:=o29+o34;
o10:=i13+o8;
o6:=o44+o8;
o3:=i6+o6;
o24:=i7+o10;
o0:=i3+o24;
o46:=i1+o0;
o28:=i2+o3;
o40:=o29+o28;
o26:=o44+o40;
o1:=o29+o0;
o27:=i9+o40;
o17:=i0+o27;
o11:=i4+o17;
o23:=i5+o11;
o4:=o44+o11;
o37:=i9+o4;
o45:=i13+o1;
t38:=i12+o36;
o9:=o4+t38;
o7:=o17+o9;
t41:=o3+o24;
o39:=o7+t41;
o38:=o17+t41;
o5:=i13;
o13:=i10;
o20:=i12;
o21:=i14;
o31:=i9;
o35:=i11;
\end{verbatim}
\noindent $R$: 49 XORs
\begin{verbatim}
o29:=i4+i6;
o15:=i8+i10;
o39:=i9+i10;
o16:=i0+i12;
o43:=i14+o15;
o22:=i12+o43;
o14:=i15+o22;
o3:=i11+o14;
o30:=i8+o3;
o28:=i7+o3;
o6:=i4+o3;
o19:=i0+o30;
o25:=i3+o19;
o36:=i2+o25;
o13:=i10+o36;
o7:=i11+o13;
o9:=o39+o7;
o40:=i7+o7;
o44:=i6+o40;
t19:=i8+o43;
o34:=i4+t19;
o12:=o16+t19;
t22:=o15+o14;
o38:=o39+t22;
t24:=i5+o38;
o46:=i4+t24;
o1:=t22+o46;
o5:=i6+o1;
o23:=i7+t24;
o45:=i14+o1;
o0:=i13+o45;
o10:=i5+o0;
o24:=i4+o10;
o2:=i12+o24;
o33:=o16+o24;
o32:=i1+o33;
o42:=i13+o32;
o41:=i15+o42;
o17:=i3+o41;
o20:=i2+o32;
o27:=i7+o17;
o35:=i14+o32;
o11:=i5+o27;
o37:=i1+o11;
o4:=i2+o37;
o31:=i6+o4;
t46:=o42+o27;
o26:=o12+t46;
o18:=o40+t46;
o8:=i4;
o21:=i8;
\end{verbatim}
\noindent $P$: 74 XORs
\begin{verbatim}
y41:=i34+i12;
y40:=i45+y41;
y39:=i37+y40;
y38:=i26+i12+y39;
y34:=i21+i8+i16;
y33:=i43+y41+y34;
y31:=i15+y33;
y30:=i1+i39;
y29:=i35+y31+y30;
y28:=i32+y40+y29;
y27:=i20+y28;
y25:=i4+y27+i9;
y24:=i36+y25;
y22:=i25+y39+y24;
y21:=i13+y22;
y20:=i40+i18+y21;
y19:=i7+y20;
y16:=i10+i24;
y15:=i17+y19;
y14:=i27+y15;
y13:=i42+y14;
y12:=i41+y13;
y11:=i14+y12;
y10:=i5+i8+i39+y27;
y9:=i33+y16;
y8:=i31+y30+i9+y21;
y6:=i46+i38;
y5:=i23+i38;
y4:=i2+y9;
y3:=i30+y5+y4;
y2:=i22+y11;
y1:=i19+y3;
o0:=y6+y2;
o1:=i0+i24+y14;
o2:=i29+i18+y8;
o3:=i28+y19+y2+i3;
o4:=i6+y5+y1+i3;
o5:=i11+y38+y16+y15;
o6:=i44+y38+y25+y20+y10;
o7:=y38+y19+y1;
o8:=y33+y11+y6;
o9:=y29+y13;
o10:=y31+y8;
o11:=y22+y12;
o12:=y34+y4;
o13:=y28+y9;
o14:=i16+y10;
o15:=y24+y3;
\end{verbatim}
\end{multicols}


\end{document}
